\section{Conclusion}
\label{sec:conclusion}

We introduced Prospection-Guided Retrieval (PGR), a retrieval policy that is independent of the underlying memory substrate. Instead of matching only on retrospective similarity to the current turn, PGR runs an iterative \emph{simulate $\rightarrow$ retrieve $\rightarrow$ refine} loop so that simulated futures surface latent and counterfactual facts that embedding or graph search often misses. This addresses a recurring limitation in standard RAG and GraphRAG pipelines, where retrieval is largely fixed by the index structure and behaves like a narrow lookup even when users need long-horizon, goal-directed assistance. We also introduced MemoryQuest, a multi-session benchmark that requires coordinating several dated references under low query--reference embedding similarity, and we evaluated PGR on MemoryQuest, ImplexConv, and PersonaMem. Across these datasets, PGR improves recall (including full multi-reference recovery) and produces answers preferred by both LLM judges and human annotators. Overall, these results support prospection as a complementary retrieval strategy when an assistant must assemble scattered personal context across sessions rather than return the closest passage to a single query.

